% =========================================================
% Cauchy Sequences — Curated Exercise Set (40 Exercises)
% Sources: Abbott, Bartle & Sherbert, Bruckner, Tao, Lebl
% =========================================================

\section*{Cauchy Sequences — Exercise Review Set}

These exercises develop the theory of Cauchy sequences and their
relationship to convergence and completeness.

\bigskip

% ---------------------------------------------------------
% Exercise stub
% ---------------------------------------------------------

\newcommand{\cauchyex}[4]{
\subsection*{Exercise: #1}

\textbf{Source:} #2

\textbf{Technique:} #3

\textbf{Task.} #4

\begin{remark*}[Work Area]
Write your proof or computation here.
\end{remark*}

\bigskip
}

% =========================================================
% ABBOTT
% =========================================================

\section*{Exercises from Abbott}

\cauchyex
{Abbott 2.6.1}
{Abbott, \textit{Understanding Analysis}}
{tech:unpack-definition}
{Define a Cauchy sequence using the $\varepsilon$–$N$ definition.}

\cauchyex
{Abbott 2.6.2}
{Abbott}
{tech:triangle-inequality}
{Prove that every convergent sequence is Cauchy.}

\cauchyex
{Abbott 2.6.3}
{Abbott}
{tech:triangle-inequality}
{Show that if $(a_n)$ is Cauchy then it is bounded.}

\cauchyex
{Abbott 2.6.4}
{Abbott}
{tech:contradiction}
{Prove that a Cauchy sequence cannot diverge to infinity.}

\cauchyex
{Abbott 2.6.5}
{Abbott}
{tech:subsequence-extraction}
{Show that every subsequence of a Cauchy sequence is Cauchy.}

\cauchyex
{Abbott 2.6.6}
{Abbott}
{tech:triangle-inequality}
{Show that if $(a_n)$ and $(b_n)$ are Cauchy then $(a_n+b_n)$ is Cauchy.}

\cauchyex
{Abbott 2.6.7}
{Abbott}
{tech:triangle-inequality}
{Show that if $(a_n)$ is Cauchy then $(|a_n|)$ is Cauchy.}

\cauchyex
{Abbott 2.6.8}
{Abbott}
{tech:choose-N}
{Prove that if $(a_n)$ is Cauchy and has a convergent subsequence, then the whole sequence converges.}

% =========================================================
% BARTLE & SHERBERT
% =========================================================

\section*{Exercises from Bartle \& Sherbert}

\cauchyex
{Bartle 3.3.1}
{Bartle \& Sherbert}
{tech:unpack-definition}
{Define a Cauchy sequence in $\mathbb{R}$.}

\cauchyex
{Bartle 3.3.2}
{Bartle \& Sherbert}
{tech:triangle-inequality}
{Prove that convergent sequences are Cauchy.}

\cauchyex
{Bartle 3.3.3}
{Bartle \& Sherbert}
{tech:monotone-bound}
{Show that every Cauchy sequence is bounded.}

\cauchyex
{Bartle 3.3.4}
{Bartle \& Sherbert}
{tech:subsequence-extraction}
{Show that subsequences of Cauchy sequences are Cauchy.}

\cauchyex
{Bartle 3.3.5}
{Bartle \& Sherbert}
{tech:triangle-inequality}
{Show that the sum of two Cauchy sequences is Cauchy.}

\cauchyex
{Bartle 3.3.6}
{Bartle \& Sherbert}
{tech:triangle-inequality}
{Show that the product of two Cauchy sequences is Cauchy.}

\cauchyex
{Bartle 3.3.7}
{Bartle \& Sherbert}
{tech:counterexample}
{Give an example of a Cauchy sequence in $\mathbb{Q}$ that does not converge in $\mathbb{Q}$.}

\cauchyex
{Bartle 3.3.8}
{Bartle \& Sherbert}
{tech:tail-argument}
{Show that tails of Cauchy sequences remain Cauchy.}

% =========================================================
% BRUCKNER
% =========================================================

\section*{Exercises from Bruckner}

\cauchyex
{Bruckner 2.5.1}
{Bruckner, \textit{Real Analysis}}
{tech:unpack-definition}
{Define Cauchy sequences.}

\cauchyex
{Bruckner 2.5.2}
{Bruckner}
{tech:triangle-inequality}
{Prove that convergent sequences are Cauchy.}

\cauchyex
{Bruckner 2.5.3}
{Bruckner}
{tech:triangle-inequality}
{Show that Cauchy sequences are bounded.}

\cauchyex
{Bruckner 2.5.4}
{Bruckner}
{tech:subsequence-extraction}
{Show subsequences of Cauchy sequences are Cauchy.}

\cauchyex
{Bruckner 2.5.5}
{Bruckner}
{tech:triangle-inequality}
{Show sums of Cauchy sequences are Cauchy.}

\cauchyex
{Bruckner 2.5.6}
{Bruckner}
{tech:triangle-inequality}
{Show products of Cauchy sequences are Cauchy.}

\cauchyex
{Bruckner 2.5.7}
{Bruckner}
{tech:counterexample}
{Construct a Cauchy sequence in $\mathbb{Q}$ that does not converge in $\mathbb{Q}$.}

\cauchyex
{Bruckner 2.5.8}
{Bruckner}
{tech:choose-N}
{Show that a Cauchy sequence with a convergent subsequence converges.}

% =========================================================
% TAO
% =========================================================

\section*{Exercises from Tao}

\cauchyex
{Tao 6.4.1}
{Tao, \textit{Analysis I}}
{tech:unpack-definition}
{Define a Cauchy sequence.}

\cauchyex
{Tao 6.4.2}
{Tao}
{tech:triangle-inequality}
{Show that convergent sequences are Cauchy.}

\cauchyex
{Tao 6.4.3}
{Tao}
{tech:monotone-bound}
{Show that Cauchy sequences are bounded.}

\cauchyex
{Tao 6.4.4}
{Tao}
{tech:triangle-inequality}
{Show sums of Cauchy sequences are Cauchy.}

\cauchyex
{Tao 6.4.5}
{Tao}
{tech:triangle-inequality}
{Show products of Cauchy sequences are Cauchy.}

\cauchyex
{Tao 6.4.6}
{Tao}
{tech:subsequence-extraction}
{Show subsequences of Cauchy sequences are Cauchy.}

\cauchyex
{Tao 6.4.7}
{Tao}
{tech:counterexample}
{Construct a Cauchy sequence in $\mathbb{Q}$ that fails to converge.}

\cauchyex
{Tao 6.4.8}
{Tao}
{tech:choose-N}
{Show that if a Cauchy sequence has a convergent subsequence then the sequence converges.}

% =========================================================
% LEBL
% =========================================================

\section*{Exercises from Lebl}

\cauchyex
{Lebl 2.3.1}
{Lebl, \textit{Basic Analysis I}}
{tech:unpack-definition}
{Define Cauchy sequences.}

\cauchyex
{Lebl 2.3.2}
{Lebl}
{tech:triangle-inequality}
{Prove that convergent sequences are Cauchy.}

\cauchyex
{Lebl 2.3.3}
{Lebl}
{tech:monotone-bound}
{Show that Cauchy sequences are bounded.}

\cauchyex
{Lebl 2.3.4}
{Lebl}
{tech:subsequence-extraction}
{Show that subsequences of Cauchy sequences are Cauchy.}

\cauchyex
{Lebl 2.3.5}
{Lebl}
{tech:triangle-inequality}
{Show sums of Cauchy sequences are Cauchy.}

\cauchyex
{Lebl 2.3.6}
{Lebl}
{tech:triangle-inequality}
{Show products of Cauchy sequences are Cauchy.}

\cauchyex
{Lebl 2.3.7}
{Lebl}
{tech:counterexample}
{Construct a Cauchy sequence in $\mathbb{Q}$ that does not converge in $\mathbb{Q}$.}

\cauchyex
{Lebl 2.3.8}
{Lebl}
{tech:choose-N}
{Show that if a Cauchy sequence has a convergent subsequence then the sequence converges.}
